\chapter{量子动力学}
\section{时间演化算符}
\subsection{Noether定理}
每一种对称性都对应着一个守恒量，每一个守恒量都对应一种对称变换的生成元.
\begin{itemize}
    \item \textbf{空间平移对称性 $\rightarrow$ 动量 $\hat{p}$ 守恒} \\
    若空间处处均匀，物体会倾向于保持它的运动状态.
    \item \textbf{时间平移对称性 $\rightarrow$ 能量 $\hat{H}$ 守恒} \\
    若物理规律不随时间改变，系统的总能量不能凭空产生或消失.
    \item \textbf{旋转对称性 $\rightarrow$ 角动量 $\hat{L}$ 守恒} \\
    若空间没有偏向哪个方向，旋转的物体就会保持它的旋转状态.
\end{itemize}
由无穷小空间平移算符$\hat{g}(\dd{x})=1+\frac{i}{\hbar}\hat{p}\dd{x}$猜测无穷小时间演化(平移)算符$\hat{u}(\dd{t})$：
\[\hat{g}(\dd{x})\psi(x_0) = \psi(x_0+\dd{x}) \qquad \hat{u}(\dd{t})\psi(t_0) = \psi(t_0+\dd{t})\]
由于 $\ket{\psi}$ 中不包含坐标信息， $\hat{g}(\dd{x})$ 是作用在左矢上的.而 $\ket{\psi}$中 包含时间信息， $\hat{u}$ 作用在右矢上， $\hat{u}$ 与 $\hat{g}(\dd{x})$ 应是复共轭的.
\[\hat{g}(\dd{x})=1+\frac{i}{\hbar}\hat{p}\dd{x} \quad \rightarrow \quad \hat{u}(\dd{t})=1-\hat{H}\dd{t}\]
猜测时间演化算符$\hat{U}(t)=e^{-\frac{i}{\hbar}\hat{H}t}$；任意时刻形式$\hat{U}(t,t_0)=e^{-\frac{i}{\hbar}\hat{H}(t-t_0)}$.

\subsection{时间演化算符的性质}
\begin{enumerate}
    \item 非相对论量子力学中粒子不会创生和湮灭(粒子数守恒) \\
    这表明粒子总得在空间的某个或某些地方，全空间总概率必然是归一化的.所以$\hat{U}(t)$必然是幺正的，确保概率归一(保模长).
    \begin{proof}
        $\hat{U}^\dagger(t,t_0)\hat{U}(t,t_0)=\hat{I}$
        \[\begin{cases}
            \hat{U}(t,t_0) &= e^{-\frac{i}{\hbar}\hat{H}(t,t_0)} \\
            \hat{U}^\dagger(t,t_0) &= e^{\frac{i}{\hbar}\hat{H}(t,t_0)}
        \end{cases} \quad \Rightarrow \quad \hat{U}^\dagger(t,t_0)\hat{U}(t,t_0)=\hat{I}\]
    \end{proof}
    \item 从$t_0 \rightarrow t_1$再从$t_1 \rightarrow t_2$应等价于$t_0 \rightarrow t_2$ \\
    即$\hat{U}(t_2,t_0)=\hat{U}(t_2,t_1)\hat{U}(t_1,t_0)$.
    \item 从$t_0 \rightarrow t_1$应等于从$t_1 \rightarrow t_0$的逆变换 \\
    即$\hat{U}(t_0,t_1) = \hat{U}^{-1}(t_1,t_0)$.
    \item 从$t_0 \rightarrow t_0$的时间演化算符应为单位算符
    即$\hat{U}(t_0,t_0)=\hat{I}$.
\end{enumerate}

\subsection{\texorpdfstring{自旋$S=1/2$体系中的时间演化}{自旋S等于1/2体系中的时间演化}}
\begin{question}
    设一粒子一开始处于$S_x=\frac{\hbar}{2}$的本征态，现加入沿$z$轴的磁场$\vec{B}=B\vec{e_x}$.求：
    \begin{enumerate}
        \item 任意$t$时刻粒子处于什么状态？
        \item 任意$t$时刻处于$S_x=\frac{\hbar}{2}$态的概率是多少？处于$S_x=-\frac{\hbar}{2}$态的概率是多少？
        \item 任意$t$时刻$\expval{\hat{S}_x},\expval{\hat{S}_y},\expval{\hat{S}_z}$是多少？
    \end{enumerate}
\end{question}
\begin{solution}
    \begin{enumerate}
        \item $$\ket{\psi(t)} = \hat{U}(t)\ket{\psi(0)}$$
        式中$\ket{\psi(0)} = \ket{+},\hat{U}(t) = e^{-i\hat{H}t/\hbar}$. \\
        在本题中Hamilton量又称为磁偶极相互作用能(Zeeman能)，它表述为：
        \[\begin{aligned}
        \hat{H} &= -\hat{\mu}_z \cdot B \\
        &= -\gamma_s \hat{S}_z \cdot B \quad \text{式中旋磁比 } \gamma_s = \frac{e}{2\mu_c} \cdot g_s \\
        &= -\gamma_s B \hat{S}_z
        \end{aligned}\]
        显然 $\hat{U}(t)$ 是 $\hat{S}_z$ 的函数，不能直接作用在 $\ket{+}$ 态上，故而将 $\ket{+}$ 在 $\hat{S}_z$ 的本征态 $\ket{\uparrow}, \ket{\downarrow}$ 上展开：
        \[\begin{aligned}
        \ket{\psi(t)} &= e^{i\gamma_s B \hat{S}_z t / \hbar} \ket{+} \\
        &= e^{i\gamma_s B \hat{S}_z t / \hbar} \textcolor{GreenDeep}{(\ket{\uparrow}\bra{\uparrow} + \ket{\downarrow}\bra{\downarrow})} \ket{+} \quad \textcolor{GreenDeep}{\text{恒等算符}} \\
        &= e^{i\gamma_s B \hat{S}_z t / \hbar} \left( \ket{\uparrow} [1, 0] \frac{1}{\sqrt{2}} \begin{bmatrix} 1 \\ 1 \end{bmatrix} + \ket{\downarrow} [0, 1] \frac{1}{\sqrt{2}} \begin{bmatrix} 1 \\ 1 \end{bmatrix} \right) \\
        &= \frac{1}{\sqrt{2}} e^{i\gamma_s B \hat{S}_z t / \hbar} (\ket{\uparrow} + \ket{\downarrow}) \\
        &= \frac{1}{\sqrt{2}} \left( e^{i\gamma_s B \frac{\hbar}{2} t / \hbar} \ket{\uparrow} + e^{i\gamma_s B (-\frac{\hbar}{2}) t / \hbar} \ket{\downarrow} \right) \quad \text{本征方程} \\
        &= \frac{1}{\sqrt{2}} \begin{bmatrix} e^{i \gamma_s B \frac{t}{2}} \\ e^{-i \gamma_s B \frac{t}{2}} \end{bmatrix} \quad \text{令 } \omega = \frac{1}{2} \gamma_s B
        \end{aligned}\]
        则$\ket{\psi(t)} = \frac{1}{\sqrt{2}} \begin{bmatrix} e^{i\omega t} \\ e^{-i\omega t} \end{bmatrix}$.

        \vspace{1em}\hrule

        \item \[\begin{aligned}
        P_{S_x=\frac{\hbar}{2}} &= \left| \braket{+}{\psi(t)} \right|^2 \\
        &= \left| \frac{1}{\sqrt{2}} \begin{bmatrix} 1 & 1 \end{bmatrix} \frac{1}{\sqrt{2}} \begin{bmatrix} e^{i\omega t} \\ e^{-i\omega t} \end{bmatrix} \right|^2 \\
        &= \frac{1}{4} \left| e^{i\omega t} + e^{-i\omega t} \right|^2 \\
        &= \frac{1}{4} \left| \cos(\omega t) + i\sin(\omega t) + \cos(\omega t) - i\sin(\omega t) \right|^2 \\
        &= \frac{1}{4} \left| 2\cos(\omega t) \right|^2 \\
        &= \cos^2(\omega t) \\[1.5em]
        P_{S_x=-\frac{\hbar}{2}} &= 1 - P_{S_x=\frac{\hbar}{2}} \\
        &= 1 - \cos^2(\omega t) \\
        &= \sin^2(\omega t)
        \end{aligned}\]

        \hrule

        \item \[\begin{aligned}
        \ev{\hat{S}_x} &= \mel{\psi(t)}{\hat{S}_x}{\psi(t)} \\
        &= \frac{1}{\sqrt{2}} \begin{bmatrix} e^{-i\omega t} & e^{i\omega t} \end{bmatrix} \frac{\hbar}{2} \begin{bmatrix} 0 & 1 \\ 1 & 0 \end{bmatrix} \frac{1}{\sqrt{2}} \begin{bmatrix} e^{i\omega t} \\ e^{-i\omega t} \end{bmatrix} \\
        &= \frac{\hbar}{4} (-e^{-i2\omega t} + e^{i2\omega t}) \\
        &= \frac{\hbar}{2} \cos(2\omega t) \\[1.5em]
        \ev{\hat{S}_y} &= \mel{\psi(t)}{\hat{S}_y}{\psi(t)} \\
        &= \frac{1}{\sqrt{2}} \begin{bmatrix} e^{-i\omega t} & e^{i\omega t} \end{bmatrix} \frac{\hbar}{2} \begin{bmatrix} 0 & -i \\ i & 0 \end{bmatrix} \frac{1}{\sqrt{2}} \begin{bmatrix} e^{i\omega t} \\ e^{-i\omega t} \end{bmatrix} \\
        &= \frac{i\hbar}{4} (-e^{-i2\omega t} + e^{i2\omega t}) \\
        &= -\frac{\hbar}{2} \sin(2\omega t) \\[1.5em]
        \ev{\hat{S}_z} &= \mel{\psi(t)}{\hat{S}_z}{\psi(t)} \\
        &= \frac{1}{\sqrt{2}} \begin{bmatrix} e^{-i\omega t} & e^{i\omega t} \end{bmatrix} \frac{\hbar}{2} \begin{bmatrix} 1 & 0 \\ 0 & -1 \end{bmatrix} \frac{1}{\sqrt{2}} \begin{bmatrix} e^{i\omega t} \\ e^{-i\omega t} \end{bmatrix} \\
        &= \frac{\hbar}{4} \begin{bmatrix} e^{-i\omega t} & e^{i\omega t} \end{bmatrix} \begin{bmatrix} e^{i\omega t} \\ -e^{-i\omega t} \end{bmatrix} \\
        &= \frac{\hbar}{4} (e^0 - e^0) \\
        &= 0
        \end{aligned}\]
    \end{enumerate}
\end{solution}
上述值表明，粒子在$x-y$平面做进动.

\noindent
\begin{minipage}[c]{0.3\linewidth}
\centering
\begin{tikzpicture}[scale=1.5, extended line/.style={shorten >=-#1,shorten <=-#1}]
    \draw[->, ultra thick, dashed] (0, -0.5) -- (0, 0);
    \draw[->, ultra thick] (0, 0) -- (0, 3) node[anchor=south] {$z$};
    \draw[->, very thick] (-0.3, 2) -- (-0.3, 2.8) node[anchor=south] {$\vec{B}$};
    \coordinate (O) at (0, 0);
    \coordinate (P0) at (-1.5, 0);
    \draw[dashed, color=OrangeDeep, thick] (O) ellipse (1.5cm and 0.4cm);
    \draw[->, color=OrangeDeep, thick] (1.5, 0) arc (0:-20:1.5cm and 0.4cm);
    \draw[->, color=OrangeDeep, thick] (-1.5, 0) arc (180:160:1.5cm and 0.4cm);
    \draw[-, color=BlueDeep, ultra thick] (O) -- (P0);
    \filldraw[color=PurpleDeep!50, fill=purple!20, thick] (P0) circle (0.15cm);
    \filldraw[color=PurpleDeep, fill=white, thick] (P0) circle (0.1cm);
    \begin{scope}[shift={(P0)}, scale=0.8]
        \draw[->, color=cyan!80!blue, thick] (-0.2, -0.2) arc (-135:-45:0.25cm);
        \draw[->, color=cyan!80!blue, thick] (0.2, 0.2) arc (45:135:0.25cm);
    \end{scope}
\end{tikzpicture}
\end{minipage}
\begin{minipage}[c]{0.6\linewidth}
    \raggedright
    这是一种特殊的\textbf{Larmor}\overtext{进动}{precession}. \\
    若不仅在$x-y$平面内绕$z$轴进动，则是一般的Larmor进动.
\end{minipage}

\section{Schrödinger绘景和Heisenberg绘景}